% !TeX root = ../navier-stokes-manuscript.tex
% ============================================================
% 1.5 Class Membership and Participation
% ============================================================
\subsection{Class Membership and Participation}\label{sub:ClassMembershipAndParticipation}

The Silver Standard ends by asking for one principle that can test many different losses of smoothness against the same boundary.
That boundary is class membership.

Fix the smooth divergence-free initial datum \(u_0\), the domain prescribed by the problem, and one constant viscosity \(\nu>0\).
Let \(Q=(u,p)\) denote a proposed terminal outcome of the maximal classical solution generated by those data, together with its preterminal values.
Write \(\operatorname{Member}(Q)\) when that proposed outcome still belongs to the original admissible fluid class: it extends the distinguished preterminal solution and obeys the same incompressible fixed-viscosity law,
\[
  \partial_tu+(u\cdot\nabla)u+\nabla p=\nu\Delta u,
  \qquad
  \nabla\cdot u=0,
  \qquad
  u(0)=u_0,
\]
with pressure remaining the pressure forced by that velocity field.
The prescribed decay or periodicity and the finite-energy requirement of the original problem remain part of the same membership test.
Membership is tied to the same datum and the same preterminal solution, so an unrelated weak or distributional solution is not another candidate for \(Q\).
Smoothness is not part of this definition.
The Silver burden is to show that every nonsmooth terminal proposal fails one of the defining admissibility requirements for this original solution.

Define class exit by the Boolean complement
\[
  \operatorname{Exit}(Q):=\neg\operatorname{Member}(Q).
\]
Gold seeks to prove, for every terminal proposal extending the distinguished maximal classical solution,
\[
  \operatorname{Member}(Q)
  \Longrightarrow
  \operatorname{Smooth}(Q).
\]
Silver seeks its contrapositive
\[
  \neg\operatorname{Smooth}(Q)
  \Longrightarrow
  \operatorname{Exit}(Q).
\]
The two lines are logically equivalent because \(\operatorname{Exit}(Q)\) means \(\neg\operatorname{Member}(Q)\).
The equivalence changes the direction of the proof without completing it.
The definition of \(\operatorname{Exit}\) is automatic; proving that every nonsmooth terminal proposal exits remains the mathematical burden.

Gold and Silver now meet at one boundary.
Gold approaches it from inside the class and asks what prevents an admitted evolution from reaching nonsmoothness.
Silver begins with a proposed nonsmooth outcome and asks which requirement of the original fluid has become impossible there.
Together they turn the vague phrase \emph{not Navier--Stokes} into a concrete membership test.

\newpage
\divider{What Participation Means}

The original requirements are experienced by the fluid as one motion.
Along a particle path \(X(t)\) satisfying \(\dot X(t)=u(t,X(t))\), the particle velocity changes according to
\[
  \frac{d}{dt}u(t,X(t))
  =
  \partial_tu+(u\cdot\nabla)u
  =
  -\nabla p+\nu\Delta u.
\]
The left-hand side is the material acceleration of the same fluid whose spatial configuration determines both terms on the right.
Viscosity responds to local velocity curvature, pressure responds to the global constraint of incompressibility, and transport carries the resulting velocity through the field.
The divergence condition prevents those responses from producing an independent compression or expansion of the fluid volume.

This joined response is what we call \emph{participation}.
A point participates at a regular time when its local velocity, pressure, and derivatives satisfy the coupled momentum and incompressibility laws of the field around it.
A packet participates when the same coupled law holds throughout the packet and remains related to the surrounding field through transport, viscosity, and the global pressure closure.
Participation concerns lawful response within one field; continuity across an alleged terminal limit is a separate property to be tested rather than part of this definition.
It remains compatible with a particle being instantaneously motionless, with a derivative vanishing, or with a steady flow.

Couette flow gives the viscous part of this response a simple physical form.
Place fluid between parallel plates and let one plate move tangentially relative to the other.
The steady profile in the interior is
\[
  u(x)=(\gamma x_2,0,0),
  \qquad
  p=\text{constant},
\]
where \(\gamma\) is the shear rate.
The interior profile has \((u\cdot\nabla)u=0\) and \(\Delta u=0\), yet the viscous shear stress \(\nu\partial_{x_2}u_1=\nu\gamma\) transmits the boundary momentum continuously through the layers.
The example isolates viscous communication in a transparent setting.
In a general flow, that communication remains joined to transport, incompressibility, and the nonlocal pressure response.

Class membership can now be read physically in one direction.
Every admitted smooth continuation has local participation in the same fixed-viscosity incompressible response.
A demonstrated failure of that joined response gives a concrete class-exit witness, while continuity and finite control still require their own endpoint tests.
